% !TeX program = pdflatex
% Conversion de Questions-révision-examen1-2026.pptx.
% Source autonome, sans image ni fichier complémentaire.
% Navigation : menu, question, réponse. 14 questions et 29 diapositives.
% Les cases vides du menu sont réservées aux questions à ajouter.
\documentclass[aspectratio=43,11pt]{beamer}
\usepackage[utf8]{inputenc}
\usepackage[T1]{fontenc}
\usepackage[french]{babel}
\usepackage{lmodern,amsmath,amssymb,booktabs,array, verbatim}
\definecolor{bleulaval}{RGB}{25,62,98}
\definecolor{bleuclair}{RGB}{232,240,247}
\setbeamercolor{structure}{fg=bleulaval}
\setbeamercolor{frametitle}{fg=bleulaval,bg=bleuclair}
\setbeamercolor{block title}{fg=white,bg=bleulaval}
\setbeamercolor{block body}{fg=black,bg=bleuclair}
\setbeamercolor{button}{fg=white,bg=bleulaval}
\setbeamertemplate{navigation symbols}{}
\setbeamertemplate{footline}{%
  \hspace{0.35cm}\hyperlink{menu}{\beamergotobutton{Menu}}%
  \hfill{\scriptsize STT-2100 -- Révision de l'examen 1}\hfill
  {\scriptsize\insertframenumber/\inserttotalframenumber}\hspace{0.35cm}\vskip5pt}
\setbeamersize{text margin left=0.6cm,text margin right=0.6cm}
\hypersetup{pdftitle={STT-2100 -- Questions de révision -- Examen 1 -- 2026},
 pdfauthor={Audrey-Anne Vallée},bookmarksopen=true}
\newcommand{\entetequestion}[1]{\textbf{Vrai ou faux?}\par\medskip #1}
\newcommand{\versreponse}[1]{\vfill\hfill\hyperlink{r-#1}{\beamergotobutton{Voir la réponse}}}
\newcommand{\versquestion}[1]{\vfill\hfill\hyperlink{q-#1}{\beamergotobutton{Revenir à la question}}}
\newcommand{\petittableau}[1]{\begin{center}\fontsize{8.5}{10}\selectfont
  \renewcommand{\arraystretch}{1.15}#1\end{center}}
% Tableaux retranscrits à partir des sorties figurant dans le PowerPoint.
\newcommand{\coefficientsage}{\petittableau{%
\begin{tabular}{lrrrr}\toprule
\texttt{Coefficients:} & \texttt{Estimate} & \texttt{Std. Error} & \texttt{t value} & \texttt{Pr(>|t|)}\\\midrule
\texttt{(Intercept)} & 516113.1 & 45311.2 & 11.390 & $5.64\times10^{-13}$\\
\texttt{age} & $-1249.3$ & 804.7 & $-1.553$ & \textbf{???}\\\bottomrule
\end{tabular}}}
\newcommand{\coefficientslog}{\petittableau{%
\begin{tabular}{lrrrr}\toprule
\texttt{Coefficients:} & \texttt{Estimate} & \texttt{Std. Error} & \texttt{t value} & \texttt{Pr(>|t|)}\\\midrule
\texttt{(Intercept)} & \texttt{AAAA} & 0.31654 & \texttt{????} & $<2\times10^{-16}$\\
\texttt{I(log(income))} & $-0.51179$ & \texttt{BBBB} & $-9.992$ & $<2\times10^{-16}$\\\bottomrule
\end{tabular}}}
\begin{document}
\begin{frame}[label=menu]{Jeu questionnaire -- Examen 1}
\pdfbookmark[0]{Menu des questions}{signet-menu}
\begin{center}
\renewcommand{\arraystretch}{1.5}
\begin{tabular}{>{\centering\arraybackslash}p{2.3cm}>{\centering\arraybackslash}p{2.3cm}>{\centering\arraybackslash}p{2.3cm}>{\centering\arraybackslash}p{2.3cm}}
\toprule
{\small\bfseries\shortstack{Modèle et\\estimation}} & {\small\bfseries Inférence} & {\small\bfseries R} & {\small\bfseries Multivarié}\\\midrule
\hyperlink{q-me100}{\textbf{100}} & \hyperlink{q-in100}{\textbf{100}} & \hyperlink{q-r100}{\textbf{100}} & \hyperlink{q-m100}{\textbf{100}}\\
\hyperlink{q-me200}{\textbf{200}} & \hyperlink{q-in200}{\textbf{200}} & \hyperlink{q-r200}{\textbf{200}} & \hyperlink{q-m200}{\textbf{200}} \\
\hyperlink{q-me300}{\textbf{300}} & \hyperlink{q-in300}{\textbf{300}} & \hyperlink{q-r300}{\textbf{300}} & \hyperlink{q-m300}{\textbf{300}} \\
\hyperlink{q-me400}{\textbf{400}} & \hyperlink{q-in400}{\textbf{400}} & \hyperlink{q-r400}{\textbf{400}} & \hyperlink{q-m400}{\textbf{400}} \\
\hyperlink{q-me500}{\textbf{500}} & \hyperlink{q-in500}{\textbf{500}} & \hyperlink{q-r500}{\textbf{500}} & \hyperlink{q-m500}{\textbf{500}} \\
\hyperlink{q-me600}{\textbf{600}} & \hyperlink{q-in600}{\textbf{600}} & \hyperlink{q-r600}{\textbf{600}} & \hyperlink{q-m600}{\textbf{600}} \\
\bottomrule
\end{tabular}
\end{center}
\end{frame}
% Cases réservées : Inférence 600; R 100, 400 et 600;
% Multivarié 100, 200, 300, 400, 500 et 600.
% Pour compléter une case, ajouter sa paire de diapositives question/réponse
% et remplacer la case vide par un hyperlien vers le label de la question.

% Modèle et estimation
% Ancienne question : Théorie RLS -- 100
\begin{frame}[label=q-me100]{Modèle et estimation -- 100}
\framesubtitle{Question}
\pdfbookmark[1]{Modèle et estimation -- 100}{signet-me100}
\entetequestion{Si on ne rejette pas l'hypothèse $\beta_1=0$ dans une analyse de régression linéaire simple, on peut conclure qu'il n'y a pas de relation entre les deux variables analysées.}
\versreponse{me100}
\end{frame}

\begin{frame}[label=r-me100]{Modèle et estimation -- 100}
\framesubtitle{Réponse}
{\small \entetequestion{Si on ne rejette pas l'hypothèse $\beta_1=0$ dans une analyse de régression linéaire simple, on peut conclure qu'il n'y a pas de relation entre les deux variables analysées.}}
\begin{block}{Réponse}\textbf{Faux.} On conclut uniquement que les données ne permettent pas de détecter une relation linéaire significative. Une relation non linéaire peut notamment exister.\end{block}
\versquestion{me100}
\end{frame}

% Ancienne question : Théorie RLS -- 200
\begin{frame}[label=q-me200]{Modèle et estimation -- 200}
\framesubtitle{Question}
\pdfbookmark[1]{Modèle et estimation -- 200}{signet-me200}
Quelle est la quantité que l'on minimise pour trouver les estimateurs des moindres carrés de $\beta_0$ et $\beta_1$ en RLS?
\versreponse{me200}
\end{frame}

\begin{frame}[label=r-me200]{Modèle et estimation -- 200}
\framesubtitle{Réponse}
{\small Quelle est la quantité que l'on minimise pour trouver les estimateurs des moindres carrés de $\beta_0$ et $\beta_1$ en RLS?}
\begin{block}{Réponse}La somme des carrés des erreurs :\[Q(\beta_0,\beta_1)=\sum_{i=1}^n \varepsilon_i^2=\sum_{i=1}^n(Y_i-\beta_0-\beta_1x_i)^2.\]\end{block}
\versquestion{me200}
\end{frame}

% Ancienne question : Théorie RLS -- 300
\begin{frame}[label=q-me300]{Modèle et estimation -- 300}
\framesubtitle{Question}
\pdfbookmark[1]{Modèle et estimation -- 300}{signet-me300}
\entetequestion{La covariance entre $\beta_0$ et $\beta_1$ est toujours nulle.}
\versreponse{me300}
\end{frame}

\begin{frame}[label=r-me300]{Modèle et estimation -- 300}
\framesubtitle{Réponse}
{\small \entetequestion{La covariance entre $\beta_0$ et $\beta_1$ est toujours nulle.}}
\begin{block}{Réponse}\textbf{Vrai.} Il s'agit de la covariance entre deux constantes.\par\medskip Attention à distinguer les paramètres $\beta_0,\beta_1$ des estimateurs aléatoires $\widehat\beta_0,\widehat\beta_1$.\end{block}
\versquestion{me300}
\end{frame}

% Ancienne question : Postulats -- 400
\begin{frame}[label=q-me400]{Modèle et estimation -- 400}
\framesubtitle{Question}
\pdfbookmark[1]{Modèle et estimation -- 400}{signet-me400}
À quoi sert le postulat 4 sur la normalité des termes d'erreur?
\versreponse{me400}
\end{frame}

\begin{frame}[label=r-me400]{Modèle et estimation -- 400}
\framesubtitle{Réponse}
{\small À quoi sert le postulat 4 sur la normalité des termes d'erreur?}
\begin{block}{Réponse}Il permet de déduire les lois exactes des estimateurs,  puis celles des statistiques de test, et de construire les intervalles de confiance et des tests d'hypothèses.\end{block}
\versquestion{me400}
\end{frame}

% Ancienne question : Théorie RLS -- 500
\begin{frame}[label=q-me500]{Modèle et estimation -- 500}
\framesubtitle{Question}
\pdfbookmark[1]{Modèle et estimation -- 500}{signet-me500}
Donnez deux façons d'améliorer la précision de l'estimateur des moindres carrés de $\beta_1$, lorsque $\sigma^2$ est fixée.
\versreponse{me500}
\end{frame}

\begin{frame}[label=r-me500]{Modèle et estimation -- 500}
\framesubtitle{Réponse}
{\small Donnez deux façons d'améliorer la précision de l'estimateur des moindres carrés de $\beta_1$, lorsque $\sigma^2$ est fixée.}
\begin{block}{Réponse}\begin{itemize}\item Augmenter la taille de l'échantillon.\item Choisir des individus dont les valeurs de $x$ sont plus dispersées.\end{itemize}\[\operatorname{Var}(\widehat\beta_1)=\frac{\sigma^2}{S_{xx}},\qquad S_{xx}=\sum_i(x_i-\bar x)^2.\]\end{block}
\versquestion{me500}
\end{frame}

% Ancienne question : Postulats -- 500
\begin{frame}[label=q-me600]{Modèle et estimation -- 600}
\framesubtitle{Question}
\pdfbookmark[1]{Modèle et estimation -- 600}{signet-me600}
Si je fais l'hypothèse d'un modèle linéaire dans lequel les erreurs ont une distribution uniforme entre $-1$ et $1$ au lieu d'une distribution normale, est-ce que les estimateurs des moindres carrés sont quand même sans biais pour les coefficients de la régression? Pourquoi?
\versreponse{me600}
\end{frame}

\begin{frame}[label=r-me600]{Modèle et estimation -- 600}
\framesubtitle{Réponse}
{\small Si je fais l'hypothèse d'un modèle linéaire dans lequel les erreurs ont une distribution uniforme entre $-1$ et $1$ au lieu d'une distribution normale, est-ce que les estimateurs des moindres carrés sont quand même sans biais pour les coefficients de la régression? Pourquoi?}
\begin{block}{Réponse}\textbf{Oui.} La propriété nécessaire au calcul de l'espérance des estimateurs est $E(\varepsilon_i)=0$.\par\medskip Elle est respectée pour une loi uniforme sur $[-1,1]$.\end{block}
\versquestion{me600}
\end{frame}

% Inférence
% Ancienne question : Inférence -- 100
\begin{frame}[label=q-in100]{Inférence -- 100}
\framesubtitle{Question}
\pdfbookmark[1]{Inférence -- 100}{signet-in100}
Nommez trois façons de prendre une décision pour le test d'hypothèses\[H_0:\beta_1=0\qquad\text{contre}\qquad H_1:\beta_1\ne0.\]
\versreponse{in100}
\end{frame}

\begin{frame}[label=r-in100]{Inférence -- 100}
\framesubtitle{Réponse}
{\small Nommez trois façons de prendre une décision pour le test d'hypothèses\[H_0:\beta_1=0\qquad\text{contre}\qquad H_1:\beta_1\ne0.\]}
\begin{block}{Réponse}\begin{itemize}\item Région critique ($t_{obs}$ versus le quantile théorique).\item Valeur-p.\item Intervalle de confiance sur la pente.\end{itemize}\end{block}
\versquestion{in100}
\end{frame}

% Ancienne question : Inférence -- 200
\begin{frame}[label=q-in200]{Inférence -- 200}
\framesubtitle{Question}
\pdfbookmark[1]{Inférence -- 200}{signet-in200}
\entetequestion{On ajuste une RLS à 40 observations et les postulats sont respectés. On obtient $\widehat\beta_1=3{,}49$ et une valeur-p de $0{,}04$ pour le test bilatéral d'égalité à 0. On peut conclure que la pente est significativement différente de zéro au seuil de $5\,\%$.}
\versreponse{in200}
\end{frame}

\begin{frame}[label=r-in200]{Inférence -- 200}
\framesubtitle{Réponse}
{\small \entetequestion{On ajuste une RLS à 40 observations et les postulats sont respectés. On obtient $\widehat\beta_1=3{,}49$ et une valeur-p de $0{,}04$. On peut conclure que la pente est significativement différente de zéro au seuil de $5\,\%$.}}
\begin{block}{Réponse}\textbf{Vrai.} Puisque $0{,}04<0{,}05$, on rejette $H_0:\beta_1=0$ au seuil de $5\,\%$.\end{block}
\versquestion{in200}
\end{frame}

% Ancienne question : Inférence -- 300
\begin{frame}[label=q-in300]{Inférence -- 300}
\framesubtitle{Question}
\pdfbookmark[1]{Inférence -- 300}{signet-in300}
Nommez trois facteurs qui influencent la taille d'un intervalle de confiance pour $E(Y\mid X=x_0)$.
\versreponse{in300}
\end{frame}

\begin{frame}[label=r-in300]{Inférence -- 300}
\framesubtitle{Réponse}
{\small Nommez trois facteurs qui influencent la taille d'un intervalle de confiance pour $E(Y\mid X=x_0)$.}
\begin{block}{Réponse}Trois réponses parmi les suivantes :\begin{itemize}\item Le niveau de confiance.\item La taille de l'échantillon.\item La variance des erreurs.\item La distance entre $x_0$ et la moyenne des $x_i$.\item La dispersion des $x_i$.\end{itemize}\end{block}
\versquestion{in300}
\end{frame}

% Ancienne question : Inférence -- 400
\begin{frame}[label=q-in400]{Inférence -- 400}
\framesubtitle{Question}
\pdfbookmark[1]{Inférence -- 400}{signet-in400}
Pourquoi la marge d'erreur d'une prévision individuelle en $x^*$ est-elle différente de la marge d'erreur associée à la moyenne de $Y$ en $x^*$?
\versreponse{in400}
\end{frame}

\begin{frame}[label=r-in400]{Inférence -- 400}
\framesubtitle{Réponse}
{\small Pourquoi la marge d'erreur d'une prévision individuelle en $x^*$ est-elle différente de la marge d'erreur associée à la moyenne de $Y$ en $x^*$?}
\begin{block}{Réponse}Une observation individuelle varie autour de la droite de régression.\par\medskip Sa prévision doit tenir compte de l'incertitude sur la moyenne estimée et de la variabilité supplémentaire due au nouveau terme d'erreur. L'intervalle de prévision est donc plus large.\end{block}
\versquestion{in400}
\end{frame}

% Ancienne question : Inférence -- 500
\begin{frame}[label=q-in500]{Inférence -- 500}
\framesubtitle{Question}
\pdfbookmark[1]{Inférence -- 500}{signet-in500}
Comment calcule-t-on la statistique observée servant à tester\[H_0:\beta_0=100\qquad\text{contre}\qquad H_1:\beta_0>100?\]
\versreponse{in500}
\end{frame}

\begin{frame}[label=r-in500]{Inférence -- 500}
\framesubtitle{Réponse}
{\small Comment calcule-t-on la statistique observée servant à tester\[H_0:\beta_0=100\qquad\text{contre}\qquad H_1:\beta_0>100?\]}
\begin{block}{Réponse}On utilise la statistique observée \[T=\frac{\widehat\beta_0-100}{\sqrt{\widehat\sigma^2\left(\dfrac1n+\dfrac{\bar x^2}{S_{xx}}\right)}}.\]Sous $H_0$, $T\sim t_{n-2}$. \end{block}
\versquestion{in500}
\end{frame}

\begin{frame}[label=q-in600]{Inférence -- 600}
\framesubtitle{Question}
\pdfbookmark[1]{Inférence -- 600}{signet-in600}
Pour un jeu de données, un bon ajustement doit produire des valeurs faibles à la fois pour \(SC_{\mathrm{reg}}\) et \(SC_{\mathrm{res}}\). Vrai ou faux?
\versreponse{in600}
\end{frame}

\begin{frame}[label=r-in600]{Inférence -- 600}
\framesubtitle{Réponse}
{\small Pour un jeu de données fixé, un bon ajustement doit produire des valeurs faibles à la fois pour \(SC_{\mathrm{reg}}\) et \(SC_{\mathrm{res}}\). Vrai ou faux?}
\begin{block}{Réponse}
Faux.
\begin{itemize}
    \item \(SC_{\mathrm{tot}}\) mesure la variabilité des observations autour de leur moyenne.
\item \(SC_{\mathrm{reg}}\) mesure la part de cette variabilité expliquée par la régression.
\item \(SC_{\mathrm{res}}\) mesure la part qui reste inexpliquée.
\end{itemize}
Puisque \(SC_{\mathrm{tot}}=SC_{\mathrm{reg}}+SC_{\mathrm{res}}\), pour un total fixé, un meilleur ajustement correspond à une plus grande part expliquée et une plus petite part résiduelle.
\end{block}
\versquestion{in600}
\end{frame}

% R
\begin{frame}[fragile, label=q-r100]{R -- 100}
\framesubtitle{Question}
\pdfbookmark[1]{R -- 100}{signet-r100}
On exécute :
\begin{verbatim}
    mod <- lm(Duree ~ Distance, data = Bixi)
\end{verbatim}
Quelle est la variable réponse? Quelle est la variable explicative? Le modèle comprend-il une ordonnée à l’origine?
\versreponse{r100}
\end{frame}

\begin{frame}[fragile, label=r-r100]{R -- 100}
\framesubtitle{Réponse}
On exécute :
\begin{verbatim}
    mod <- lm(Duree ~ Distance, data = Bixi)
\end{verbatim}
Quelle est la variable réponse? Quelle est la variable explicative? Le modèle comprend-il une ordonnée à l’origine?
\begin{block}{Réponse}
\begin{itemize}
    \item Réponse: Duree
    \item Variable explicative: Distance
    \item Oui
\end{itemize}
\end{block}
\versquestion{r100}
\end{frame}

% Ancienne question : R et SAS -- 200
\begin{frame}[fragile, label=q-r200]{R -- 200}
\framesubtitle{Question}
\pdfbookmark[1]{R -- 200}{signet-r200}
Quel est le coefficient de détermination de cet ajustement?
\begin{verbatim}
> anova(mod)
Analysis of Variance Table

Response: y
          Df Sum Sq Mean Sq F value    Pr(>F)
x          1 646.26  646.26  71.492 1.638e-08 
Residuals 23 207.91    9.04
\end{verbatim}
\versreponse{r200}
\end{frame}

\begin{frame}[label=r-r200]{R -- 200}
\framesubtitle{Réponse}
{\small Quel est le coefficient de détermination de cet ajustement?\par\medskip}
\begin{block}{Réponse}\[R^2=\frac{SC_{\mathrm{reg}}}{SC_{\mathrm{tot}}}=\frac{646{,}26}{207,91+646{,}26}\simeq0{,}757.\]La régression explique environ $76\,\%$ de la variabilité observée de la variable réponse.\end{block}
\versquestion{r200}
\end{frame}

% Ancienne question : R et SAS -- 300
\begin{frame}[label=q-r300]{R -- 300}
\framesubtitle{Question}
\pdfbookmark[1]{R -- 300}{signet-r300}
On veut prédire le prix d'une maison à partir de son âge en se basant sur un échantillon de taille 35. Comment calcule-t-on la valeur-p de l'âge et quelle est sa valeur approximative?\par\medskip\coefficientsage
\versreponse{r300}
\end{frame}

\begin{frame}[label=r-r300]{R -- 300}
\framesubtitle{Réponse}
{\small On veut prédire le prix d'une maison à partir de son âge en se basant sur un échantillon de taille 35. Comment calcule-t-on la valeur-p de l'âge et quelle est sa valeur approximative?\par\medskip}
\begin{block}{Réponse}Pour le test bilatéral de pente nulle,\[p=2P\!\left(t_{33}> |-1{,}553|\right)\simeq 0{,}130.\]Les degrés de liberté sont $n-2 = 35-2=33$.\end{block}
\versquestion{r300}
\end{frame}

\begin{frame}[fragile, label=q-r400]{R -- 400}
\framesubtitle{Question}
\pdfbookmark[1]{R -- 400}{signet-r400}
\begin{verbatim}
mod <- lm(Duree ~ Distance, data = Bixi)
nouveau <- data.frame(distance = 1.8)
predict(mod, nouveau, interval = "confidence", level = 0.90)
predict(mod, nouveau, interval = "prediction", level = 0.90)
\end{verbatim}
Que représente chaque intervalle? Leurs centres sont-ils identiques? Lequel est le plus large?
\versreponse{r400}
\end{frame}

\begin{frame}[fragile, label=r-r400]{R -- 400}
\framesubtitle{Réponse}
\begin{verbatim}
mod <- lm(Duree ~ Distance, data = Bixi)
nouveau <- data.frame(distance = 1.8)
predict(mod, nouveau, interval = "confidence", level = 0.90)
predict(mod, nouveau, interval = "prediction", level = 0.90)
\end{verbatim}
Que représente chaque intervalle? Leurs centres sont-ils identiques? Lequel est le plus large?
\begin{block}{Réponse}
Le premier concerne la durée moyenne des déplacements de 1,8 km, le second un nouveau déplacement individuel de 1,8 km. Leurs centres sont identiques, mais l’intervalle de prévision est plus large.
\end{block}
\versquestion{r400}
\end{frame}


% Ancienne question : R et SAS -- 500
\begin{frame}[label=q-r500]{R -- 500}
\framesubtitle{Question}
\pdfbookmark[1]{R -- 500}{signet-r500}
Donnez la valeur de \texttt{BBBB}.\par\bigskip\coefficientslog
\versreponse{r500}
\end{frame}

\begin{frame}[label=r-r500]{R -- 500}
\framesubtitle{Réponse}
{\small Donnez la valeur de \texttt{BBBB}.\par\bigskip}
\begin{block}{Réponse}La statistique du test d'une pente nulle est\[t_{\mathrm{obs}}=\frac{\widehat\beta_1}{\operatorname{se}(\widehat\beta_1)}=\frac{-0{,}51179}{\texttt{BBBB}}.\]Ainsi,\[\texttt{BBBB}=\frac{-0{,}51179}{-9{,}992}\simeq0{,}0512.\]\end{block}
\versquestion{r500}
\end{frame}

\begin{frame}[label=q-r600]{R -- 600}
\framesubtitle{Question}
\pdfbookmark[1]{R -- 600}{signet-r600}
De mémoire, est-ce que la sortie de \texttt{summary(modele)} nous donne directement la valeur de l'estimateur de $\sigma^2$?
\versreponse{r600}
\end{frame}

\begin{frame}[label=r-r600]{R -- 600}
\framesubtitle{Réponse}
{\small De mémoire, est-ce que la sortie de \texttt{summary(modele)} nous donne directement la valeur de l'estimateur de $\sigma^2$?}
\begin{block}{Réponse}
Non! Elle nous donne la valeur de l'estimateur de $\sigma$.
\end{block}
\versquestion{r600}
\end{frame}

%Multivarié

\begin{frame}[label=q-m100]{Multivarié -- 100}
\framesubtitle{Question}
\pdfbookmark[1]{Multivarié -- 100}{signet-m100}
Soit une fonction scalaire $f: \mathbb{R}^n \rightarrow \mathbb{R}$ et un vecteur $\bold x = \left( x_1, \dots, x_n\right)^\top$. Le résultat de la dérivée partielle $\displaystyle \frac{\partial f(\bold x)}{ \partial \bold x}$ est un scalaire. Vrai ou faux? 
\versreponse{m100}
\end{frame}

\begin{frame}[label=r-m100]{Multivarié -- 100}
\framesubtitle{Réponse}
{\small Soit une fonction scalaire $f: \mathbb{R}^n \rightarrow \mathbb{R}$ et un vecteur $\bold x = \left( x_1, \dots, x_n\right)^\top$. Le résultat de la dérivée partielle $\displaystyle \frac{\partial f(\bold x)}{ \partial \bold x}$ est un scalaire. Vrai ou faux? }
\begin{block}{Réponse}
Faux. On obtient un vecteur de longueur $n$, \begin{align*}
\frac{\partial f(\bold x)}{ \partial \bold x} =
    \begin{pmatrix}
       \displaystyle \frac{\partial f(\bold x)}{\partial x_1} \\[4mm]
       \displaystyle \frac{\partial f(\bold x)}{\partial x_2} \\[4mm]
       \displaystyle \vdots\\[4mm]
       \displaystyle \frac{\partial f(\bold x)}{\partial x_n} 
    \end{pmatrix}.
\end{align*}
\end{block}
\versquestion{m100}
\end{frame}

\begin{frame}[label=q-m200]{Multivarié -- 200}
\framesubtitle{Question}
\pdfbookmark[1]{Multivarié -- 200}{signet-m200}
On projette un vecteur \(y\) à l’aide d’une matrice de projection orthogonale \(P\), puis on projette à nouveau le résultat sur le même espace. Le résultat change-t-il? Quelle propriété permet de répondre?
\versreponse{m200}
\end{frame}

\begin{frame}[label=r-m200]{Multivarié -- 200}
\framesubtitle{Réponse}
{\small On projette un vecteur \(y\) à l’aide d’une matrice de projection orthogonale \(P\), puis on projette à nouveau le résultat sur le même espace. Le résultat change-t-il? Quelle propriété permet de répondre?}
\begin{block}{Réponse}
Non. Puisque \(P\) est idempotente,
\[
P(Py)=P^2y=Py.
\]
Le vecteur projeté appartient déjà à l’espace sur lequel on projette.
\end{block}
\versquestion{m200}
\end{frame}

\begin{frame}[label=q-m300]{Multivarié -- 300}
\framesubtitle{Question}
\pdfbookmark[1]{Multivarié -- 300}{signet-m300}
Soit un vecteur aléatoire $\mathbf X=(X_1,X_2)^\top$ qui suit une loi normale multivariée, ainsi $\mathbf X \sim \mathcal{N}_2 \left( \boldsymbol{\mu}_X, \quad \boldsymbol{\Sigma}_X \right)$. On a que
\begin{align*}
\boldsymbol\mu_X=\begin{pmatrix}2\\3\end{pmatrix},\qquad
\boldsymbol\Sigma_X=\begin{pmatrix}4&3\\3&6\end{pmatrix}.
\end{align*}
Les composantes de $\mathbf X$ sont-elles homoscédastiques? Indépendantes? 
\versreponse{m300}
\end{frame}

\begin{frame}[label=r-m300]{Multivarié -- 300}
\framesubtitle{Réponse}
{\small Soit un vecteur aléatoire $\mathbf X=(X_1,X_2)^\top$ qui suit une loi normale multivariée, ainsi $\mathbf X \sim \mathcal{N}_2 \left( \boldsymbol{\mu}_X, \quad \boldsymbol{\Sigma}_X \right)$. On a que
\begin{align*}
\boldsymbol\mu_X=\begin{pmatrix}2\\3\end{pmatrix},\qquad
\boldsymbol\Sigma_X=\begin{pmatrix}4&3\\3&6\end{pmatrix}.
\end{align*}
Les composantes de $\mathbf X$ sont-elles homoscédastiques? Indépendantes? }
\begin{block}{Réponse}
Les composantes ne sont pas homoscédastiques puisque leurs variances diffèrent. $\mathrm{Var}(X_1) = 4 \neq 6 = \mathrm{Var}(X_2)$. Elles ne sont pas indépendantes puisque $\operatorname{Cov}(X_1,X_2)=3\ne0$.
\end{block}
\versquestion{m300}
\end{frame}

\begin{frame}[label=q-m400]{Multivarié -- 400}
\framesubtitle{Question}
\pdfbookmark[1]{Multivarié -- 400}{signet-m400}
Trois variables aléatoires normales indépendantes ont une même moyenne \(\mu\), mais des variances respectives \(1\), \(4\) et \(9\). Écrivez le vecteur des espérances et la matrice de variances-covariances.
\versreponse{m400}
\end{frame}

\begin{frame}[label=r-m400]{Multivarié -- 400}
\framesubtitle{Réponse}
{\small Trois variables aléatoires normales indépendantes ont une même moyenne \(\mu\), mais des variances respectives \(1\), \(4\) et \(9\). Écrivez le vecteur des espérances et la matrice de variances-covariances.}
\begin{block}{Réponse}
\[
\boldsymbol\mu=
\begin{pmatrix}\mu\\\mu\\\mu\end{pmatrix},
\qquad
\Sigma=
\begin{pmatrix}
1&0&0\\
0&4&0\\
0&0&9
\end{pmatrix}.
\]
\end{block}
\versquestion{m400}
\end{frame}

\begin{frame}[label=q-m500]{Multivarié -- 500}
\framesubtitle{Question}
\pdfbookmark[1]{Multivarié -- 500}{signet-m500}
On a
\[
X=\begin{pmatrix}X_1\\X_2\end{pmatrix}
\sim N_2\left(
\begin{pmatrix}2\\-1\end{pmatrix},
\begin{pmatrix}4&0\\0&9\end{pmatrix}
\right).
\]
Pour \(T=X_1+2X_2\), calculez \(E(T)\) et \(\operatorname{Var}(T)\). 
\versreponse{m500}
\end{frame}

\begin{frame}[label=r-m500]{Multivarié -- 500}
\framesubtitle{Réponse}
{\small On a
\[
X=\begin{pmatrix}X_1\\X_2\end{pmatrix}
\sim N_2\left(
\begin{pmatrix}2\\-1\end{pmatrix},
\begin{pmatrix}4&0\\0&9\end{pmatrix}
\right).
\]
Pour \(T=X_1+2X_2\), calculez \(E(T)\) et \(\operatorname{Var}(T)\). }
\begin{block}{Réponse}
\[
E(T)=2+2(-1)=0,
\]
\[
\operatorname{Var}(T)
=4+4(9) + 2(0)=40.
\]
\end{block}
\versquestion{m500}
\end{frame}

\begin{frame}[label=q-m600]{Multivarié -- 600}
\framesubtitle{Question}
\pdfbookmark[1]{Multivarié -- 600}{signet-m600}
Soit  $\mathbf X \sim \mathcal{N}_2 \left( \boldsymbol{\mu}_X, \quad \boldsymbol{\Sigma}_X \right)$, avec
\begin{align*}
\boldsymbol\Sigma_X=\begin{pmatrix}4&3\\3&6\end{pmatrix}.
\end{align*}
Soit
\[ \mathbf Y=\mathbf C\mathbf X,\qquad
\mathbf C=\begin{pmatrix}1&0\\0&1\end{pmatrix}.\]
Calculez et donnez les valeurs de $\operatorname{Var}(\mathbf Y)$. 
\versreponse{m600}
\end{frame}

\begin{frame}[label=r-m600]{Multivarié -- 600}
\framesubtitle{Réponse}
{\small Soit  $\mathbf X \sim \mathcal{N}_2 \left( \boldsymbol{\mu}_X, \quad \boldsymbol{\Sigma}_X \right)$, avec
\begin{align*}
\boldsymbol\Sigma_X=\begin{pmatrix}4&3\\3&6\end{pmatrix}.
\end{align*}
Soit
\[ \mathbf Y=\mathbf C\mathbf X,\qquad
\mathbf C=\begin{pmatrix}1&0\\0&1\end{pmatrix}.\]
Calculez et donnez les valeurs de $\operatorname{Var}(\mathbf Y)$. }
\begin{block}{Réponse}
\[\operatorname{Var}(\mathbf Y)=\mathbf C\boldsymbol\Sigma_X\mathbf C^\top
=\begin{pmatrix}1&0\\0&1\end{pmatrix}\begin{pmatrix}4&3\\3&6\end{pmatrix}\begin{pmatrix}1&0\\0&1\end{pmatrix}
=\begin{pmatrix}4&3\\3&6\end{pmatrix}.\]
\end{block}
\versquestion{m200}
\end{frame}


\end{document}





